El presente capítulo establece la terminología oficial y el significado de las siglas utilizadas a lo largo de este Manual de Gestión de Seguridad Operacional (SMS). Su objetivo es garantizar una interpretación uniforme, precisa y estandarizada de los conceptos por parte de todo el personal de SKYFILMS, la Aerocivil y terceros involucrados en la operación.\\

\textbf{1. Acrónimos y Siglas}

Para los fines del presente manual, se aplicarán las siguientes siglas y acrónimos:

\begin{longtable}{|>{\centering\arraybackslash}m{0.17\linewidth}|m{0.6\linewidth}|}
\hline
\rowcolor{gray!20}
\textbf{Sigla / Acrónimo} & \textbf{Significado (Inglés / Español)} \\
\hline
\endfirsthead % Indica que este encabezado solo se muestra en la primera página de la tabla

\hline
\rowcolor{gray!20}
\textbf{Sigla / Acrónimo} & \textbf{Significado (Inglés / Español)} \\
\hline
\endhead % Indica que este encabezado se repetirá automáticamente si la tabla pasa a otra hoja

AOC / CDO & Certificado de Explotador de Servicios Aéreos / Certificado de Operador. \\ \hline
C2 & Command and Control / Enlace de Mando y Control del UAS. \\ \hline
ERP / PRE & Emergency Response Plan / Plan de Respuesta ante Emergencias. \\ \hline
FRAT & Flight Risk Assessment Tool / Herramienta de Evaluación de Riesgos de Vuelo. \\ \hline
KPI & Key Performance Indicator / Indicador Clave de Rendimiento. \\ \hline
MAC & Mid-Air Collision / Colisión en el Aire. \\ \hline
MIE & Matriz de Identificación y Evaluación de Riesgos. \\ \hline
MOR & Mandatory Occurrence Report / Reporte Obligatorio de Eventos de Seguridad Operacional. \\ \hline
OACI / ICAO & Organización de Aviación Civil Internacional / International Civil Aviation Organization. \\ \hline
RAC & Reglamentos Aeronáuticos de Colombia. \\ \hline
RPA & Remotely Piloted Aircraft / Aeronave Pilotada a Distancia. \\ \hline
SA & Safety Assurance / Aseguramiento de la Seguridad Operacional. \\ \hline
SMS & Safety Management System / Sistema de Gestión de Seguridad Operacional. \\ \hline
SPI & Safety Performance Indicator / Indicador de Rendimiento de Seguridad Operacional. \\ \hline
SPT & Safety Performance Target / Meta de Rendimiento de Seguridad Operacional. \\ \hline
SRM & Safety Risk Management / Gestión de Riesgos de Seguridad Operacional. \\ \hline
TRE & Tecnología, Reglamento y Entrenamiento (Metodología de mitigación de riesgos RAC 100). \\ \hline
Aerocivil & Unidad Administrativa Especial de Aeronáutica Civil. \\ \hline
UAS & Unmanned Aircraft System / Sistema de Aeronave No Tripulada. \\ \hline
VLOS & Visual Line of Sight / Operación en Línea de Vista Visual. \\ \hline
DIACC & Grupo de Investigación de Accidentes de la Aerocivil. \\ \hline
FFHH & Factores Humanos. \\ \hline
SAG & Safety Action Group / Grupo de Acción de Seguridad Operacional. \\ \hline
SRB & Safety Review Board / Junta de Revisión de Seguridad Operacional. \\ \hline
SDCPS & Safety Data Collection and Processing Systems / Sistemas de Recopilación y Procesamiento de Datos de Seguridad Operacional. \\ \hline
VOR & Voluntary Occurrence Report / Reporte Voluntario. \\ \hline
DIS & Documento de Información Suplementaria entregado a la Aerocivil. \\ \hline
OpSpecs & Especificaciones de Operación. \\ \hline
MO & Manual de Operaciones de Skyfilms. \\ \hline
MCM & Manual de Control de Mantenimiento de Skyfilms. \\ \hline
RUAS & Registro de Unidades de Aeronaves no Tripuladas. \\ \hline
FPV & First Person View / Vista en Primera Persona. \\ \hline
EPP & Elementos de Protección Personal. \\ \hline
SST & Seguridad y Salud en el Trabajo. \\ \hline

% Nota: En longtable la leyenda (caption) y etiqueta (label) se ponen adentro, al final de la tabla
\caption{Acrónimos del Manual SMS}
\label{tab:Acronimos}
\end{longtable}

\textbf{2. Definiciones}

Con el fin de asegurar una comprensión integral de los procesos de SKYFILMS, se adoptan las siguientes definiciones oficiales (alineadas con el RAC 1 y RAC 219):

\begin{enumerate}
    \item \textbf{Aceptabilidad del Riesgo:} El nivel de riesgo que una organización está dispuesta a asumir tras haber implementado las defensas y controles correspondientes, basándose en la matriz de tolerabilidad.

    \item \textbf{Auditoría de Seguridad Operacional:} Examen sistemático, independiente y documentado para determinar si las actividades y los resultados del SMS cumplen con las disposiciones planificadas y la normatividad aeronáutica.

    \item \textbf{Consecuencia:} El resultado potencial, impacto o efecto de que un peligro se materialice (Ej. Daños a terceros, lesiones, pérdida del equipo).

    \item \textbf{Cultura Justa (Just Culture):} Un entorno donde no se castiga a los empleados por acciones, omisiones o decisiones tomadas de acuerdo con su experiencia y capacitación, pero en el cual no se toleran la negligencia grave, las infracciones intencionales ni los actos destructivos.

    \item \textbf{Ejecutivo Responsable:} Persona única en la organización (SKYFILMS) que actúa al más alto nivel directivo y cuyas responsabilidades, basadas en su autoridad y atribuciones financieras, son absoluta y totalmente indelegables frente a la eficacia del SMS.

    \item \textbf{Gestión del Cambio:} Proceso formal para identificar, evaluar y mitigar los riesgos de seguridad operacional que surgen como resultado de modificaciones en la estructura de la empresa, adquisición de nuevos equipos (UAS), o cambios normativos.

    \item \textbf{Mejora Continua:} Proceso recurrente que permite incrementar la eficacia y eficiencia del Sistema de Gestión de Seguridad Operacional mediante la evaluación de indicadores (SPIs), auditorías y lecciones aprendidas.

    \item \textbf{Mitigación (Defensas):} Medida, control o acción correctiva adoptada para reducir la probabilidad o la severidad de los riesgos asociados a un peligro identificado. En SKYFILMS, se categorizan bajo el modelo TRE (Tecnología, Reglamento, Entrenamiento).

    \item \textbf{Peligro:} Condición, objeto o actividad que podría provocar un incidente o accidente de aviación, o contribuir al mismo.

    \item \textbf{Probabilidad:} La posibilidad de que ocurra un evento inseguro o una consecuencia indeseada al interactuar con un peligro.

    \item \textbf{Rendimiento en Materia de Seguridad Operacional:} El logro de una organización con respecto a la seguridad operacional, definido mediante la medición de sus indicadores (SPI) y metas (SPT).

    \item \textbf{Riesgo de Seguridad Operacional:} La evaluación combinada de la probabilidad y la severidad proyectadas de las consecuencias de un peligro.

    \item \textbf{Seguridad Operacional (Safety):} El estado en el cual los riesgos asociados a las actividades de aviación (operación de UAS) se reducen y controlan a un nivel aceptable.

    \item \textbf{Severidad:} La magnitud o gravedad del impacto de las consecuencias de un peligro en caso de que este se materialice.

    \item \textbf{Sistema de Aeronave No Tripulada (UAS):} Una aeronave y sus elementos conexos, que se operan sin piloto a bordo (incluye la estación de control, el enlace C2 y los equipos de apoyo).

    \item \textbf{Gerente SMS:} Persona designada por el Gerente General para administrar el SMS. Es el único nombre de este cargo en todo el Manual (equivale a Asesor SMS o Gerente de Seguridad Operacional en otros documentos).

    \item \textbf{Junta de Revisión de Seguridad Operacional (SRB):} Grupo directivo, presidido por el Gerente General, que revisa el desempeño del SMS y decide sobre recursos y riesgos de nivel alto.

    \item \textbf{Grupo de Acción de Seguridad Operacional (SAG):} Grupo operativo, dirigido por el Gerente SMS, que analiza reportes, riesgos y acciones y las lleva a cabo.

    \item \textbf{Riesgo Residual:} Riesgo que queda después de aplicar los controles.

    \item \textbf{Riesgo Agregado:} Suma de varios riesgos pequeños que, juntos, pueden ser mayores que cada uno por separado.
\end{enumerate}